\section*{Bab \ref{ch:g2:measure} --- Uang dan Pengukuran}

\begin{solution}{exo:g2:measure:1}
Pensil: cm. Panjang ruang kelas: m. Bulu: g. Tas sekolah: kg. Botol
air: L.
\end{solution}

\begin{solution}{exo:g2:measure:2}
$2$ m $= 200$ cm; \quad $300$ cm $= 3$ m; \quad $1$ kg $= 1000$ g.
\end{solution}

\begin{solution}{exo:g2:measure:3}
Misalnya: $16 = 10 + 5 + 1$; \ $23 = 20 + 2 + 1$; \
$35 = 20 + 10 + 5$. (Cara lain juga mungkin.)
\end{solution}

\begin{solution}{exo:g2:measure:4}
$10 - 6 = 4$; \quad $20 - 17 = 3$; \quad $50 - 32 = 18$.
\end{solution}

\begin{solution}{exo:g2:measure:5}
Jarum panjang di $6$, jarum pendek antara $4$ dan $5$: pukul
setengah lima ($4{:}30$). Jarum panjang di $3$, jarum pendek tepat
setelah $8$: pukul delapan lewat seperempat ($8{:}15$).
\end{solution}

\begin{solution}{exo:g2:measure:6}
Pukul setengah tiga: jarum panjang menunjuk ke bawah ke angka $6$,
jarum pendek berada di antara angka $2$ dan $3$.
\end{solution}

\begin{solution}{exo:g2:measure:7}
Sesudah Mei: Juni. Sebelum Januari: Desember. (Bulan ulang tahun beragam.)
\end{solution}

\begin{solution}{exo:g2:measure:8}
Dibelanjakan: $3 + 4 = 7$. Kembalian: $10 - 7 = 3$.
\end{solution}

\begin{solution}{exo:g2:measure:9}
$1$ kg $= 1000$ g, yang lebih dari $800$ g: gulanya lebih berat,
\omterm{ex:g1:subtraction:difference}{selisihnya} $1000 - 800 = 200$ g.
\end{solution}

\begin{solution}{exo:g2:measure:10}
Misalnya $20 + 10 + 5 + 5 = 40$ dan $10 + 10 + 10 + 10 = 40$
(dua cara dengan tepat empat uang logam atau uang kertas).

\end{solution}
